
\documentclass{article} % For LaTeX2e
\usepackage{iclr2026_conference,times}

% Optional math commands from https://github.com/goodfeli/dlbook_notation.
\input{math_commands.tex}

\usepackage{hyperref}
\usepackage{url}
\usepackage{booktabs}
\usepackage{graphicx}

\title{Merge-Split Dynamic Tokenization for Flexibility}

\author{Justin Chae\thanks{Equal contribution.} \& Hanbyul Kang\footnotemark[1] \& Alvin Zeng\footnotemark[1] \\
Paul G.\ Allen School of Computer Science \& Engineering \\
University of Washington \\
\texttt{\{jchae3, kahaby, wnzeng\}@cs.washington.edu}
}

\newcommand{\fix}{\marginpar{FIX}}
\newcommand{\new}{\marginpar{NEW}}

\iclrfinalcopy % Uncomment for camera-ready version, but NOT for submission.
\begin{document}

\maketitle

\begin{abstract}
TODO
\end{abstract}

% ── 1. Introduction ────────────────────────────────────────────────────────────
\section{Introduction}
\label{sec:intro}

TODO

% ── 2. Related Work ────────────────────────────────────────────────────────────
\section{Related Work}
\label{sec:related}

\subsection{Subword Tokenization}
TODO

\subsection{Dynamic and Adaptive Tokenization}
TODO

\subsection{Korean NLP and Multilingual Benchmarks}
TODO

% ── 3. Method ──────────────────────────────────────────────────────────────────
\section{Method}
\label{sec:method}

\subsection{Dynamic BPE Merge}
TODO

\subsection{SuperBPE: Cross-Word Merging}
TODO

\subsection{Entropy-Guided Split}
TODO

\subsection{Merge-Split Pipeline}
TODO

\subsection{Hypernetwork Embedding}
TODO

% ── 4. Korean Tokenizer Analysis ───────────────────────────────────────────────
\section{Korean Tokenizer Analysis}
\label{sec:tokenizer}

\subsection{Effect of Korean Training Data on Tokenization Efficiency}
TODO

\subsection{Token-Level Visualization}
TODO

% ── 5. Experiments ─────────────────────────────────────────────────────────────
\section{Experiments}
\label{sec:experiments}

\subsection{Experimental Setup}
TODO

\subsection{English Benchmark: MMLU}
TODO

\subsection{Korean Benchmarks: KMMLU, KoBEST, HRM8K}
TODO

% ── 6. Results and Analysis ────────────────────────────────────────────────────
\section{Results and Analysis}
\label{sec:results}

\subsection{Merge vs.\ Split vs.\ Merge-Split}
TODO

\subsection{Entropy Threshold Sensitivity}
TODO

\subsection{English--Korean Accuracy Gap}
TODO

% ── 7. Conclusion ──────────────────────────────────────────────────────────────
\section{Conclusion}
\label{sec:conclusion}

TODO

% ── Acknowledgments ────────────────────────────────────────────────────────────
\subsubsection*{Acknowledgments}
TODO

\bibliography{iclr2026_conference}
\bibliographystyle{iclr2026_conference}

\appendix
\section{Appendix}
TODO

\end{document}

%% ============================================================
%% TEMPLATE REFERENCE (commented out — not compiled)
%% ============================================================
%
% \section{Submission of conference papers to ICLR 2026}
%
% ICLR requires electronic submissions, processed by
% \url{https://openreview.net/}. See ICLR's website for more instructions.
%
% If your paper is ultimately accepted, the statement {\tt
%   {\textbackslash}iclrfinalcopy} should be inserted to adjust the
% format to the camera ready requirements.
%
% \section{General formatting instructions}
% \label{gen_inst}
%
% The text must be confined within a rectangle 5.5~inches (33~picas) wide and
% 9~inches (54~picas) long. The left margin is 1.5~inch (9~picas).
% Use 10~point type with a vertical spacing of 11~points. Times New Roman is the
% preferred typeface throughout. Paragraphs are separated by 1/2~line space,
% with no indentation.
%
% There will be a strict upper limit of \textbf{9 pages} for the main text of
% the initial submission, with unlimited additional pages for citations.
% This limit will be expanded to \textbf{10 pages} for rebuttal/camera ready.
%
% \section{Headings: first level}
% First level headings are in small caps, flush left and in point size 12.
%
% \subsection{Headings: second level}
% Second level headings are in small caps, flush left and in point size 10.
%
% \subsubsection{Headings: third level}
% Third level headings are in small caps, flush left and in point size 10.
%
% \section{Citations, figures, tables, references}
%
% \subsection{Citations within the text}
% Use \verb|\citet{}| when authors are part of the sentence,
% and \verb|\citep{}| for parenthetical citations.
%
% \subsection{Footnotes}
% Indicate footnotes with a number in the text.
% Precede the footnote with a horizontal rule of 2~inches (12~picas).
%
% \subsection{Figures}
% \begin{figure}[h]
% \begin{center}
% \fbox{\rule[-.5cm]{0cm}{4cm} \rule[-.5cm]{4cm}{0cm}}
% \end{center}
% \caption{Sample figure caption.}
% \end{figure}
%
% \subsection{Tables}
% \begin{table}[t]
% \caption{Sample table title}
% \label{sample-table}
% \begin{center}
% \begin{tabular}{ll}
% \multicolumn{1}{c}{\bf PART}  &\multicolumn{1}{c}{\bf DESCRIPTION} \\ \hline \\
% Dendrite  & Input terminal \\
% Axon      & Output terminal \\
% Soma      & Cell body (contains cell nucleus) \\
% \end{tabular}
% \end{center}
% \end{table}
%
% \section{Default Notation}
% In an attempt to encourage standardized notation, we have included the
% notation file from the textbook, \textit{Deep Learning} \cite{goodfellow2016deep}.
%
% \centerline{\bf Numbers and Arrays}
% $\displaystyle a$ -- scalar,\quad $\displaystyle \va$ -- vector,\quad
% $\displaystyle \mA$ -- matrix,\quad $\displaystyle \tA$ -- tensor
%
% \section{Final instructions}
% Do not change any aspects of the formatting parameters in the style files.
% Please note that pages should be numbered.
%
% \section{Preparing PDF files}
% Please prepare PDF files with paper size ``US Letter''.
% Consider directly generating PDF files using \verb+pdflatex+.
%
% \begin{verbatim}
% dvips mypaper.dvi -t letter -Ppdf -G0 -o mypaper.ps
% ps2pdf mypaper.ps mypaper.pdf
% \end{verbatim}
%
% \subsubsection*{Author Contributions}
% If you'd like to, you may include a section for author contributions.
%
% \subsubsection*{Acknowledgments}
% All acknowledgments, including those to funding agencies, go at the end of the paper.
